% !TeX root=../main.tex
\chapter{نتایج شبیه‌سازی و ارزیابی عملکرد}
\label{chap:results}

\section{محیط شبیه‌سازی}
بسته‌های نرم‌افزاری، سخت‌افزار اجرا، و چارچوب تولید داده و نویز در این بخش توصیف می‌شوند.

\section{پارامترهای شبیه‌سازی}
پارامترهای پایه شامل حافظهٔ \(m\)، طول \(L\)، تعداد فریم \(F\)، احتمال خطای بیت \(P_b\)، نوع درهم‌ریز، نوع کد، خاتمه، پانچینگ و نوع مشاهدهٔ سخت/نرم است. نمونه‌های موردی مانند کدهای \((13,15)\)، \((15,17)\) و \((133,171)\) و طول‌های \(L\in\{50,100,300,500\}\) در ارزیابی لحاظ می‌شوند.

\section{معیارهای ارزیابی}
\subsection{Reconstruction Success Rate}
\subsection{False Reconstruction Rate}
\subsection{Number of admissible solutions}
\subsection{Structural distance}
\subsection{Runtime}
\subsection{Number of expanded nodes}
\subsection{Number of pruned nodes}
این معیارها فراتر از دقت صرف، جنبهٔ ابهام و هزینهٔ محاسباتی را نیز پوشش می‌دهند.

\section{اثر SNR / \(P_b\)}
منحنی‌های موفقیت و پیچیدگی برحسب سطح نویز گزارش می‌شوند.

\section{اثر طول interleaver}
تأثیر طول جایگزشت بر نرخ موفقیت و زمان اجرا بررسی می‌شود.

\section{اثر حافظه کدگذار}
عملکرد برای مقادیر مختلف \(m\) مقایسه می‌گردد.

\section{اثر تعداد فریم‌های دریافتی}
افزایش \(F\) معمولاً آمار قیود را پایدارتر می‌کند؛ این بخش حساسیت را می‌سنجد.

\section{اثر termination}
نقش مدل خاتمه در صحت طول فریم و بازسازی ساختار تحلیل می‌شود.

\section{اثر نوع interleaver}
درهم‌ریزهای تصادفی و ساخت‌یافته از نظر دشواری بازیابی مقایسه می‌شوند.

\section{اثر hard و soft observation}
تفاوت مشاهدهٔ سخت و نرم در دقت و هزینه بررسی می‌شود.

\section{تحلیل ambiguity}
حتی در حالت بدون نویز ممکن است چند بازسازی معادل وجود داشته باشد؛ این بخش توزیع تعداد جواب‌های مجاز را تحلیل می‌کند.

\section{مقایسه با روش‌های موجود}
مقایسه در چهار محور انجام می‌شود: گسترهٔ بازسازی، دانش پیشین موردنیاز، مقاومت به نویز، و پیچیدگی محاسباتی.

\section{تحلیل آماری Monte Carlo}
تنظیمات تکرار مونت‌کارلو، فاصلهٔ اطمینان و آزمون‌های آماری مرتبط گزارش می‌شوند.

\section{بحث نتایج}
یافته‌های کلیدی، محدودیت‌های مشاهده‌شده و دلالت‌های آن‌ها برای طراحی الگوریتم جمع‌بندی می‌گردند.
